\documentclass[12pt, a4paper, oneside]{article}

% Pacotes essenciais
\usepackage[utf8]{inputenc}
\usepackage[portuguese]{babel}
\usepackage{amsmath, amsfonts, amssymb}
\usepackage{graphicx}
\usepackage{geometry}
\usepackage{hyperref}
\usepackage{cite}
\usepackage{setspace}

\geometry{a4paper, left=3cm, top=3cm, right=2cm, bottom=2cm}
\onehalfspacing

\title{Desenvolvimento e Avaliação de um Repositório Guiado para Apoio ao Ensino de Programação para Pessoas com Deficiência Visual}
\author{Nome do Aluno \\ Orientadora: Nome da Orientadora}
\date{\today}

\begin{document}

\maketitle

\begin{abstract}
A inclusão de Pessoas com Deficiência Visual (PcDV) no campo da Ciência da Computação apresenta desafios significativos devido à histórica dependência de metáforas visuais em linguagens e ambientes de programação. Apesar da existência de estudos, ferramentas e diretrizes voltadas para a acessibilidade no ensino de programação, esses recursos encontram-se dispersos na literatura e na web. O objetivo deste estudo é desenvolver e avaliar um repositório centralizado, denominado \textit{IncluiDev}, que organiza e disponibiliza essas ferramentas, métodos e artigos científicos. O repositório foi construído como uma aplicação web acessível e incorpora um "Assistente Guia", um sistema de recomendação dinâmico baseado em perfis de usuário (Estudantes, Professores, Desenvolvedores e Pesquisadores). A validação do sistema será conduzida utilizando as escalas \textit{System Usability Scale} (SUS) e o \textit{Technology Acceptance Model} (TAM), juntamente com avaliações qualitativas de acessibilidade. Os resultados preliminares indicam que a centralização e a recomendação guiada reduzem o esforço de busca e promovem o acesso equitativo à educação tecnológica.
\\
\textbf{Palavras-chave:} deficiência visual; ensino de programação; acessibilidade; repositório educacional; usabilidade.
\end{abstract}

\section{Introdução}

\subsection{Contextualização}
A inclusão digital de Pessoas com Deficiência Visual (PcDV) nas áreas de Ciência da Computação e Tecnologia da Informação (TI) tem se mostrado cada vez mais essencial em uma sociedade fortemente impulsionada pela tecnologia. O ensino de programação tradicionalmente exige a assimilação de metáforas visuais complexas, como diagramas de fluxo, código em blocos coloridos e a própria indentação textual, criando barreiras significativas de aprendizagem para estudantes cegos ou com baixa visão.

Felizmente, a comunidade acadêmica e de desenvolvimento tem respondido a essas barreiras através da criação de linguagens baseadas em evidências (como a Quorum Language), ambientes tangíveis (como o Code Jumper) e a adaptação de leitores de tela em IDEs de mercado. No entanto, o simples desenvolvimento dessas soluções não garante o acesso ao usuário final.

\subsection{Problema de Pesquisa}
Apesar do crescente número de estudos sobre o ensino de programação para pessoas com deficiência visual, os recursos encontram-se profundamente dispersos, dificultando sua localização e utilização prática por diferentes públicos. Professores em busca de metodologias de ensino, desenvolvedores buscando diretrizes de acessibilidade e os próprios estudantes enfrentam a exaustiva tarefa de garimpar soluções fragmentadas. Nesse contexto, emerge a questão de pesquisa: como centralizar, classificar e recomendar de forma acessível metodologias, ferramentas e pesquisas sobre o ensino de programação para PcDV?

\subsection{Objetivos}
O objetivo geral deste trabalho é construir e avaliar um portal web dinâmico e centralizado que funcione como um repositório interativo de práticas, ferramentas e estudos voltados ao ensino de programação para pessoas com deficiência visual.

Para alcançar este propósito, definiram-se os seguintes objetivos específicos:
\begin{itemize}
    \item Organizar os recursos identificados na literatura acadêmica e na indústria em uma taxonomia coerente;
    \item Desenvolver um repositório web que respeite rigorosamente as diretrizes de acessibilidade (WCAG 2.1 AA);
    \item Implementar um sistema de recomendações ("Busca Guiada") baseado em jornadas de perfis (Estudante, Professor, Pesquisador e Desenvolvedor);
    \item Avaliar a usabilidade, a aceitação e a acessibilidade do sistema desenvolvido utilizando métodos empíricos e escalas padronizadas.
\end{itemize}

\subsection{Organização do artigo}
O restante deste artigo está organizado da seguinte forma: a Seção 2 apresenta a Fundamentação Teórica sobre o ensino acessível e os métodos de avaliação. A Seção 3 detalha os Métodos aplicados na construção do repositório e os procedimentos de validação. A Seção 4 reporta os Resultados alcançados, enquanto a Seção 5 discute as implicações desses achados. Finalmente, a Seção 6 traz a Conclusão e os trabalhos futuros.

% TODO: Preencher as seções seguintes.
\section{Fundamentação Teórica}
\textit{Em desenvolvimento...}

\section{Método}
\textit{Em desenvolvimento...}

\section{Resultados}
\textit{Em desenvolvimento...}

\section{Discussão}
\textit{Em desenvolvimento...}

\section{Conclusão}
\textit{Em desenvolvimento...}

\end{document}
